% Introduction to the rook vs bishop chapter, for the 1500-1599 book only.
% The words and theory it needs come from the library (lib/).

\section*{What to look for}
Rook against bishop is often a fight between a stronger piece and a better position. The rook
wins when it can tie the enemy pieces down; the bishop holds when its side finds a fortress, or
when the rook can be given up at the right moment. Most mistakes in this chapter came from a few
ideas.

\begin{itemize}
\item The rook's strongest weapons are \term{zugzwang}\index{Zugzwang} and the threat of a
\term{back-rank-mate}.\index{Rook!Back-rank mate} With the enemy pieces
\termas{tie-down}{tied down},\index{Strategy!Tying down} the defender soon runs out of moves
(\xref{Diagram}{depaula-miguel-56b}, \xref{Diagram}{kwiek-kozlinski-62w} and
\xref{Diagram}{butzke-silva-52w}). The same threats can also help the defender
(\xref{Diagram}{kose-ulusoy-57b}).

\item A passive rook does not win. Keep it active, so that it stops the enemy pawns and supports
your own (\term{active-pieces},\index{Strategy!Active pieces} \xref{Diagram}{chongchoryuen-lam-51w}
and \xref{Diagram}{piriz-barandiaran-45b}).

\item Count the material that remains after an exchange. Giving up the rook for the bishop and a
pawn often leaves a pawn ending that is simply won or simply drawn
(\term{transposing},\index{Transposing into a pawn ending} \xref{Diagram}{headlong-denning-56w},
\xref{Diagram}{szucs-borhy-62w} and \xref{Diagram}{dvalishvili-sakovich-61w}).

\item Many of these endings are pawn races, and every tempo counts
(\term{do-not-lose-tempo}).\index{Pawn races!Do not lose a tempo} Do not make a move that achieves
nothing (\xref{Diagram}{batzoli-louka-68b}, \xref{Diagram}{duru-atabey-54b} and
\xref{Diagram}{saglamer-acar-63w}).

\item The defender has his own resources. A \term{barrier}\index{Barrier} or a
\term{fortress}\index{Fortress} holds (\xref{Diagram}{krupacz-frac-62b} and
\xref{Diagram}{phillips-bhushan-67b}). The \term{wrong-bishop}\index{Theoretical draws!Wrong bishop and rook's pawn against king}
and the \term{ponziani}\index{Theoretical draws!Ponziani position} are draws, sometimes reached
like a miracle (\xref{Diagram}{tsarouchas-ieromonachos-68w}, \xref{Diagram}{sagir-bulut-47b} and
\xref{Diagram}{coimbra-oliveira-55w}). A rook checking from too close fails: the
\term{three-file-rule}\index{Rook!Three-file rule} tells you how far away it must be
(\xref{Diagram}{deleon-manteiga-66b}).

\item With rook and rook's pawn against bishop, the defending king must not go to the
\term{dangerous-corner}\index{Rook vs bishop!Dangerous corner} (\xref{Diagram}{georgiev-simeonov-49b}
and \xref{Diagram}{lenaerts-fuglestein-41w}).

\item Colours matter. With the bishop, put your pawns on the
\termas{pawns-other-colour}{other colour than your bishop},\index{Bishop!Pawns on the other colour than the bishop}
otherwise the enemy king walks in on the squares your bishop cannot guard
(\xref{Diagram}{molinanavarro-munozhoyos-38b}). Against the bishop, keep your pieces off its
colour:\index{Practical play!Pieces off the colour of the enemy bishop} then it can never hit them,
and you do not need to check for tactics at all (\xref{Diagram}{krupacz-frac-62b} and
\xref{Diagram}{alexakis-stefanatos-38b}).

\item Many of these endings tempt you into long calculation. Choose a plan first, and calculate
only what is forced: long variation, wrong
variation\index{Practical play!Long variation, wrong variation} (\xref{Diagram}{monchosi-costabel-43w}).
\end{itemize}
